\section{What Emergence Means, and What It Does Not Mean}

\begin{definition}[Emergence]
\textbf{Emergence} is the appearance of a higher-level organized pattern whose
behaviour cannot be usefully described as a mere list of isolated local events,
because the pattern feeds back on later local behaviour, alters future costs,
or becomes an explanatory layer in its own right.
\end{definition}

\begin{definition}[Framework formation]
\textbf{Framework formation} is the process by which repeated inputs,
conditioning, environmental coupling, and historical reuse generate a
relatively stable behavioural or cognitive regime.
\end{definition}

\begin{remark}
The chapter uses ``emergence'' in a disciplined explanatory sense, not as a
magic word for anything complicated. A pattern is not emergent merely because
it is surprising, difficult, or emotionally important.
\end{remark}

\subsection{What Emergence Does Not Mean}

For this manuscript, emergence does \emph{not} mean:
\begin{enumerate}[nosep]
  \item \textbf{Mystery inflation:} calling a process emergent because its
        micro-details are inconvenient to track.
  \item \textbf{Causal vagueness:} pretending that once a macro-pattern
        appears, local mechanisms no longer matter.
  \item \textbf{Unrestricted analogy:} assuming that because emergence occurs
        in one biological or physical system, the same mechanism must govern a
        particular human routine one-to-one.
  \item \textbf{Strong proof by vocabulary:} upgrading system language such as
        ``basin,'' ``phase,'' or ``conditioning'' into direct evidence about
        daily agency without an empirical bridge.
\end{enumerate}

\begin{discussion}
This negative definition is strategically important. Without it, the chapter
would invite exactly the criticism it should avoid: that emergence is being
used as a prestige label rather than as a constrained explanatory tool. The
book's claim is narrower. Repeated local interaction can generate
history-sensitive organization; that organization can later structure what is
easy, salient, and likely; and this makes the higher-level regime worth naming.
\end{discussion}
